% generated by figures/make_figures.py — do not edit
\begin{longtable}{lp{14.5cm}}
\toprule
Table & What it shows \\
\midrule
\endhead
\multicolumn{2}{l}{\emph{Design, data and timeline}} \\
S1 & Which labelled AOIs entered the analysis, by inclusion criterion, flood type and continent. \\
S2 & The flooded area of the excluded AOIs, by criterion. \\
S3 & The time of every step of the study, including each amendment and addendum. \\
S4 & The content of each amendment. \\
S5 & The planning status of each group of reported estimates. \\
\multicolumn{2}{l}{\emph{Pre-specified analyses}} \\
S6--S7 & The within-area comparison judged against five measures of the baseline gain and two bootstraps. \\
S8 & The primary comparison within descriptive strata (H5). \\
\multicolumn{2}{l}{\emph{Exploratory checks of the verdict}} \\
S9 & The primary comparison with one definition of permanent water, other subsampling seeds and a bare-earth elevation model. \\
S10 & The skill of flood maps, with one definition of the HAND-only map, per AOI and pooled. \\
S11 & The primary comparison by sensor, product generation and without Africa, and with the alternative second stages and references of amendment 10. \\
S12 & Brier scores of every split. \\
S13 & The activation split in two code versions, further intervals, the GFM domain and the denominator of the labels. \\
S14 & A census of the exploratory interval estimates. \\
S15 & Whether areas with more water-side names flood more, within language regions. \\
S16 & Test inversion and the other analyses of amendment 5. \\
\multicolumn{2}{l}{\emph{What assignment through settlements can carry}} \\
S17 & Gains from effects injected into the observed labels, base models fixed. \\
S18 & Bounds, detectable effects and margin translations with different no-effect references, with Monte Carlo intervals. \\
S19 & Bounds counted three ways, in both directions and for effects that differ between settlements. \\
S20 & Simulated power of the decision rule with every model refitted. \\
S21 & How an effect on the ground a name refers to is diluted on all ground nearest the name. \\
S22 & The margin and the bounds on the scale of the estimator of Sect. 5. \\
\multicolumn{2}{l}{\emph{Flooding near settlements}} \\
S23 & Element groups of the lexicon. \\
S24 & The group of island and water-meadow names at the settlement sites and without the meadow terms. \\
S25 & The pooled odds ratio near settlements under every variant, by river size and by class, and the two label sets on the same cells. \\
S26 & Where the flooded cells near water-side names lie, by country and by activation. \\
S27 & The floodplain reading by country, without one activation and without doubtful forms. \\
\multicolumn{2}{l}{\emph{How the names were read}} \\
S28 & Forms excluded by the lexicon's own rules or prone to transferred names, and the most frequent forms by country. \\
S29 & The settlements of an example AOI and their readings. \\
S30 & The fixed prompt of the language-model reading. \\
S31 & The classes of the language-model reading. \\
S32 & Cells whose names are read in a language other than their own. \\
S33 & The Wikipedia headings taken as explanations of the name, by language. \\
\multicolumn{2}{l}{\textbf{Additional tables}} \\
\multicolumn{2}{l}{\emph{Additional: Pre-specified analyses}} \\
S34 & The secondary comparisons of the protocol under each split. \\
S35 & Quantities planned in amendment 2 and not shown elsewhere, including the protocol's two-way bootstrap. \\
S36 & Ratio-bootstrap intervals of $R$ for every comparison and split. \\
S37 & Operating characteristics of the decision rule on synthetic labels for each split. \\
S38 & Secondary metrics of the held-out predictions. \\
S39 & The confirmatory comparisons with the frozen code and with the corrected reading of Norwegian names. \\
S40 & The coefficients of the second stage. \\
S41 & Interval coverage in the null simulations, the $p$ values of the protocol and the spatial correlation of residuals. \\
\multicolumn{2}{l}{\emph{Additional: Exploratory checks of the verdict}} \\
S42 & The models of the exploratory analyses. \\
S43 & The within-area comparison under exploratory variants of sample, labels and reading. \\
S44 & The primary comparison in subsets of cells and under other fold assignments. \\
S45 & The achieved equivalence bound, the location groups needed to settle the inconclusive comparisons, the permutation test, the interval with the second stage refitted and the operating characteristics at the margin. \\
S46 & The Murphy decomposition, the comparison within AOIs and the skill of flood maps. \\
S47 & What the regional composition of names adds. \\
S48 & Exploratory runs and controls of amendment 6. \\
\multicolumn{2}{l}{\emph{Additional: What assignment through settlements can carry}} \\
S49 & The same gains with every model refitted on the injected labels. \\
S50 & The estimator of Sect. 5 on labels with a known injected odds ratio. \\
S51 & Translations of the margin and bounds under the synthetic labels of the protocol. \\
S52 & Recurrent surface water around settlements, by name class (no flood labels). \\
\multicolumn{2}{l}{\emph{Additional: Flooding near settlements}} \\
S53--S54 & Where settlements with names of each class stand relative to the drainage, by distance (no flood labels). \\
S55 & Heights of settlements above the main rivers, by name class (no flood labels). \\
S56 & Within-AOI odds ratios of flooding by class and distance. \\
S57 & Odds ratios near settlements by class and distance band, with pairs-bootstrap intervals. \\
S58 & The names increment by class of the nearest name. \\
S59 & Flood contrasts and increments by distance from the settlement. \\
S60 & The pooled odds ratio near settlements with each activation left out. \\
S61 & Adjusted $p$ and $q$ values of the class coefficients. \\
S62 & The odds ratio near water-side names corrected over a grid of misreading rates. \\
S63 & The odds ratio near water-side names on every evaluable cell, with both label sets. \\
S64 & Every quantity of the analyses on GFM labels (amendment 3). \\
S65 & Diagnostics of the GFM labels. \\
S66 & Calibration and skill of the held-out predictions, by split. \\
\multicolumn{2}{l}{\emph{Additional: How the names were read}} \\
S67 & The lexicon by language block. \\
S68 & The water-side forms that matched the most cells. \\
S69 & Examples of names read as water-side, with their recorded derivations. \\
S70 & Precision and recall of the readings in more detail. \\
S71 & Recorded etymologies by lexicon class and language. \\
S72 & Agreement of a second coder and the precision and recall from its codes. \\
S73 & The settlements by GeoNames feature code, and the pilot simulations. \\
S74 & Settlements whose namesake is recorded on Wikidata. \\
S75 & Precision and recall by language. \\
S76 & Precision and recall by country and by form, for both coders. \\
\multicolumn{2}{l}{\emph{Additional: Further quantities}} \\
S77 & Quantities stated in the paper and not given in another table. \\
\bottomrule
\end{longtable}
